% Generated draft: requires mathematical and source review.
\documentclass[11pt]{article}
\usepackage[margin=25mm]{geometry}
\usepackage{fontspec}
\setmainfont{Latin Modern Roman}
\usepackage{amsmath,amssymb,mathtools,unicode-math}
\setmathfont{Latin Modern Math}
\usepackage{enumitem,needspace,xurl,hyperref}
\setlength{\parindent}{0pt}
\setlength{\parskip}{5pt}
\setlength{\emergencystretch}{3em}
\newcommand{\N}{\mathbb{N}}
\newcommand{\Z}{\mathbb{Z}}
\newcommand{\Q}{\mathbb{Q}}
\newcommand{\R}{\mathbb{R}}
\newcommand{\C}{\mathbb{C}}
\newcommand{\F}{\mathbb{F}}
\newcommand{\A}{\mathbb{A}}
\newcommand{\PP}{\mathbb{P}}
\newcommand{\E}{\mathbb{E}}
\newcommand{\eps}{\varepsilon}
\newcommand{\dd}{\,\mathrm{d}}
\begin{document}
\section*{UnsolvedMath ID: 30006244}
\textbf{Source problem code:} catalogue entry
\subsection{Definitions and mathematical statement}
Let \(p\) be a prime and let \(S\) be a finite \(p\)-group. A fusion system \(\mathcal F\) on \(S\) is a category whose objects are the subgroups of \(S\), whose morphisms are injective group homomorphisms, which contains all homomorphisms induced by conjugation in \(S\), and in which every morphism factors as an isomorphism in the category followed by an inclusion.

Write \(\operatorname{Aut}_{\mathcal F}(P)=\operatorname{Hom}_{\mathcal F}(P,P)\), and let \(\operatorname{Aut}_S(P)\) consist of automorphisms induced by \(N_S(P)\). A subgroup \(P\leq S\) is fully normalized, respectively fully centralized, if \(|N_S(P)|\), respectively \(|C_S(P)|\), is maximal among the corresponding orders for subgroups isomorphic to \(P\) in \(\mathcal F\). The fusion system is saturated if the following hold:
\begin{enumerate}
\item Every fully normalized \(P\leq S\) is fully centralized and satisfies \(\operatorname{Aut}_S(P)\in\operatorname{Syl}_p(\operatorname{Aut}_{\mathcal F}(P))\).
\item If \(\varphi\in\operatorname{Hom}_{\mathcal F}(P,S)\) and \(\varphi(P)\) is fully centralized, then \(\varphi\) extends to a morphism \(N_\varphi\to S\) in \(\mathcal F\), where, writing \(c_s(x)=sxs^{-1}\),
\[
N_\varphi=\{s\in N_S(P):\varphi\circ(c_s|_P)\circ\varphi^{-1}\in\operatorname{Aut}_S(\varphi(P))\}.
\]
\end{enumerate}

For a finite group \(G\) containing \(S\), define \(\mathcal F_S(G)\) by
\[
\operatorname{Hom}_{\mathcal F_S(G)}(P,Q)
=\{c_g|_P:P\to Q\mid g\in G,\ gPg^{-1}\leq Q\}
\qquad(P,Q\leq S).
\]
The system \(\mathcal F\) is realizable if there is such a finite group \(G\) with \(S\in\operatorname{Syl}_p(G)\) and \(\mathcal F=\mathcal F_S(G)\).

\textbf{Recoverable characterization problem.} For arbitrary primes \(p\), finite \(p\)-groups \(S\), and saturated fusion systems \(\mathcal F\) on \(S\), determine necessary and sufficient conditions for realizability:
\[
\exists\,G\text{ finite}\quad
S\leq G,\qquad S\in\operatorname{Syl}_p(G),\qquad
\mathcal F=\mathcal F_S(G).
\]

\textbf{Unresolved normal-subgroup qualification.} The supplied statement additionally requires \(H\trianglelefteq G\). It neither specifies \(H\) nor quantifies it nor describes any associated data to be preserved. If \(H\) is simply existentially quantified with no further requirements, this condition imposes no restriction, since one can take \(H=1\). If \(H\) is prescribed, the requested extension problem cannot be recovered precisely from this record. The displayed characterization problem therefore records only the unambiguous realizability part; the additional requirement involving \(H\) remains unresolved.

\textbf{Tameness qualification.} The source also asks for the relation between realizability and tameness in full generality and records the historical question whether every realizable fusion system is tame. No definition of tameness or exact remaining assertion is supplied. The supplied dated assessment attributes an affirmative result about realizable systems being tame to Broto--M{\o}ller--Oliver--Ruiz, while describing the broader characterization problem as partially solved. Consequently the historical question is retained as source context, without asserting that it remains open or specifying a new unresolved tameness conjecture.
\subsection{Short English statement}
Characterize the saturated fusion systems arising from conjugation in finite groups with the given Sylow \(p\)-subgroup. The source also requires a normal subgroup \(H\), but does not specify its role, and asks about tameness without defining the intended remaining target. Its partially solved label and reported results are supplied dated curation requiring mathematical review.
\subsection{Sources}
\begin{itemize}
\item Supplied local UnsolvedMath registry record, permanent ID \(30006244\), title \emph{Realizability and Tameness of Fusion Systems}, source problem code \texttt{catalogue entry}, in the Oberwolfach Reports Open Problems collection. The record gives proposed year 2025 and does not name a proposer.
\item \emph{Finite Groups, Fusion Systems and Applications} (2025), Oberwolfach Reports, pp.\ 689--752, DOI \texttt{10.4171/owr/2025/16}: \url{https://doi.org/10.4171/owr/2025/16}. This is the supplied source report and the literature-source link labeled ``current workshop context.'' The record states that the quoted passage is retained verbatim as context for the extracted problem.
\item Broto--M{\o}ller--Oliver--Ruiz, 2023, as attributed in the supplied literature sources: \url{https://doi.org/10.1112/plms.12571}. No article title or precise theorem reference is supplied.
\item Literature assessment dated \texttt{2026-08-22}: the supplied curation labels the problem partially solved and attributes to Broto--M{\o}ller--Oliver--Ruiz both realizability from realizable components and the assertion that realizable fusion systems are tame, while stating that a full characterization remains open. This is supplied dated curation, not an independently verified assessment. The record expressly requests independent verification of time-sensitive status and cited claims.
\end{itemize}
\end{document}
